\documentclass[11pt]{article}

\usepackage[T1]{fontenc}
\usepackage{amsmath,amssymb,amsthm}
\usepackage{hyperref}

\title{Constrained and Multiple-Group Closed-Form CFA:\\
       Minimum-Distance Projection, Measurement Invariance,\\
       and a Reference-Group Scalar Map}
\author{}
\date{\today}

\newcommand{\T}{^{\mathsf T}}
\newcommand{\E}{\mathbb{E}}
\newcommand{\Var}{\operatorname{Var}}
\newcommand{\Cov}{\operatorname{Cov}}
\newcommand{\tr}{\operatorname{tr}}
\newcommand{\diag}{\operatorname{diag}}
\newcommand{\vech}{\operatorname{vech}}
\newcommand{\vecop}{\operatorname{vec}}
\newcommand{\rank}{\operatorname{rank}}
\newcommand{\Om}{\Omega}

\theoremstyle{plain}
\newtheorem{proposition}{Proposition}
\newtheorem{corollary}{Corollary}
\theoremstyle{remark}
\newtheorem*{remark}{Remark}

\begin{document}
\maketitle

\section*{Purpose}

Guttman's (1952) multiple-group method gives a closed-form CFA estimator
$\hat\theta = \tau(s)$ as a smooth map from the sample covariance
$s = \vech(S)$ to a complete parameter vector, with delta-method covariance
$\Om = J\Gamma J\T / N$ for $J = \partial\tau/\partial s\T$ and the fourth-moment
matrix $\Gamma$ (see \emph{guttman\_cfa\_asymptotics.tex} and
\emph{noniterative\_cfa\_tests.tex}). This note adds three things while keeping
the estimator closed form: independent multiple-group blocks, general linear
equality constraints (measurement invariance and the rest) through a one-step
minimum-distance projection, and true scalar invariance through a
reference-group mean map. The projection gives an exact $\chi^2$ constraint
test. The scalar map is the one piece that is not a pure projection, because it
depends on the loadings, so its test is a linearized Wald.

\section{Multiple groups}

For groups $b = 1,\dots,G$ with independent samples, fit each block by the
same map: $\hat\theta_b = \tau(s_b)$, $J_b = \partial\tau/\partial s_b\T$,
$\Gamma_b$, $N_b$. Stack $\hat\theta = (\hat\theta_1,\dots,\hat\theta_G)$. Since
the blocks are independent, the delta-method covariance is block diagonal,
\begin{equation}
  \Om = \bigoplus_{b=1}^{G} \frac{1}{N_b}\, J_b\, \Gamma_b\, J_b\T ,
  \label{eq:omega-blockdiag}
\end{equation}
and the residual goodness-of-fit statistic and its weighted-$\chi^2$ reference
are additive across blocks: $T = \sum_b N_b\, r_b\T V_b\, r_b$ with
$r_b = \vech(S_b) - \vech(\Sigma_b(\hat\theta_b))$, joint
$\mathrm{df} = \sum_b (p_b^\ast - q_b)$, and reference spectrum equal to the
union of the per-block spectra of $(M_b\T V_b M_b)\Gamma_b$ where
$M_b = I - \Delta_b J_b$ is the residual projector. The configural fit imposes
no cross-group ties. All cross-group hypotheses are imposed afterward on
\eqref{eq:omega-blockdiag}.

\section{Linear equality constraints as a projection}

Let the constraint be $R\theta = c$ with $R$ of full row rank $k$ (the rows of
$R$ and the vector $c$ are exactly the $A_{\mathrm{eq}}$, $b_{\mathrm{eq}}$ of the
partable's \texttt{group.equal} family). The $\Om^{-1}$-metric projection of
$\hat\theta$ onto the constraint surface is
\begin{equation}
  \tilde\theta
   = \hat\theta - \Om R\T (R\Om R\T)^{-1}(R\hat\theta - c),
  \qquad
  \tilde\Om
   = \Om - \Om R\T (R\Om R\T)^{-1} R\, \Om ,
  \label{eq:proj}
\end{equation}
and $\tilde\Om$ has rank $q - k$. This is the closed-form
measurement-invariance estimator. The associated statistic is
\begin{equation}
  W = (R\hat\theta - c)\T (R\Om R\T)^{-1}(R\hat\theta - c).
  \label{eq:wald}
\end{equation}

\begin{proposition}[Wald equals minimum distance, and is exact $\chi^2_k$]
\label{prop:duality}
With $\tilde\theta$ as in \eqref{eq:proj},
\[
  W
  = (\hat\theta - \tilde\theta)\T \Om^{-1} (\hat\theta - \tilde\theta) .
\]
If $\sqrt{N}(\hat\theta - \theta_0) \to \mathcal N(0, N\Om)$ and $R\theta_0 = c$,
then $W \to \chi^2_k$ exactly, not a weighted mixture of $\chi^2_1$ variables.
\end{proposition}

\begin{proof}
Write $g = R\hat\theta - c$ and $A = R\Om R\T$. From \eqref{eq:proj},
$\hat\theta - \tilde\theta = \Om R\T A^{-1} g$, so
\[
  (\hat\theta-\tilde\theta)\T \Om^{-1}(\hat\theta-\tilde\theta)
  = g\T A^{-1} R\, \Om\, \Om^{-1}\, \Om R\T A^{-1} g
  = g\T A^{-1}(R\Om R\T)A^{-1} g
  = g\T A^{-1} g = W .
\]
For the distribution, $g = R\hat\theta - c = R(\hat\theta - \theta_0)$ under the
null, so $\sqrt N\, g \to \mathcal N(0, N R\Om R\T) = \mathcal N(0, N A)$. Hence
$g\T A^{-1} g \to \chi^2_k$. The mixture form that appears for the residual
goodness-of-fit statistic does not appear here because \eqref{eq:wald} uses the
true covariance $A$ of $g$ as its own weight, not a model-implied surrogate.
\end{proof}

\begin{remark}
The rank drop $\rank(\tilde\Om) = q-k$ is immediate: $\Om R\T A^{-1} R\,\Om$ is
$\Om$-symmetric idempotent of rank $k$, so \eqref{eq:proj} removes exactly the
$k$ constrained directions. When $c = 0$ (a pure cross-group equality) the
projection reduces to the classic constrained-GLS form. The general path keeps
$c$ so that fixed-value and effect-coding constraints are handled by the same
line of code.
\end{remark}

\paragraph{Reach.} Every hypothesis whose restriction is linear in $\theta$ is
covered by \eqref{eq:proj}--\eqref{eq:wald} in closed form with an exact
$\chi^2_k$ test. This includes metric invariance (equal loadings), strict
invariance (equal residual variances), tau-equivalence and parallelism, effect
coding, fixed values, and intercept equality with latent means fixed at zero
(Section~\ref{sec:means}). Table~\ref{tab:reach} places these against the two
harder regimes.

\begin{table}[h]
\centering
\begin{tabular}{lll}
\hline
Restriction & Form & Machinery \\
\hline
Metric / strict / $\tau$-equiv. / effect / fixed & $R\theta = c$ (linear)
  & projection \eqref{eq:proj}, exact $\chi^2_k$ \\
True scalar (free latent means) & $(I-P_\Lambda)(m_g - m_r) = 0$ (bilinear)
  & reference map, linearized Wald \\
Nonlinear equality & $h(\theta)=0$ & linearize-and-iterate (out of scope) \\
Inequality / Heywood bounds & $R\theta \ge c$ & active set (out of scope) \\
\hline
\end{tabular}
\caption{How far the closed form reaches. The first row is a pure projection.
The second is the one bilinear case, handled in Section~\ref{sec:scalar}. The
last two leave closed form.}
\label{tab:reach}
\end{table}

\section{Mean structure}
\label{sec:means}

Add free intercepts $\nu_g$ with latent means fixed at zero. The saturated
configural intercept is $\nu_g = m_g$, the sample mean, because
$\mu_g(\theta) = \nu_g + \Lambda_g\alpha_g = \nu_g$ when $\alpha_g = 0$. Order
each block's moments mean-first, $z_b = (m_b\T, \vech(S_b)\T)\T$. The estimator
Jacobian gains an analytic mean block, $\partial\nu_g/\partial m_g = I$ and
$\partial\nu_g/\partial\vech(S_g) = 0$, and the normal-theory
$\Gamma$ becomes block diagonal, $\Gamma^{+}_{\mathrm{NT}} =
\diag(\Sigma_b, \Gamma_{\mathrm{NT}}(\Sigma_b))$, because normal-theory third
moments vanish.

\begin{proposition}[The mean part is inert for fit]
\label{prop:mean-inert}
Order parameters intercepts-first. The augmented residual projector
$M^{+} = I - \Delta^{+} J^{+}$ over $z_b$ is block diagonal with a zero
intercept block:
\[
  \Delta^{+} J^{+}
  = \begin{pmatrix} I & 0 \\ 0 & \Delta_b J_b \end{pmatrix},
  \qquad
  M^{+} = \begin{pmatrix} 0 & 0 \\ 0 & M_b \end{pmatrix}.
\]
Hence the goodness-of-fit spectrum of $(M^{+\mathsf T} V^{+} M^{+})\Gamma^{+}$
equals the covariance-only spectrum of $(M_b\T V_b M_b)\Gamma_b$ together with
zeros, and the joint degrees of freedom are unchanged: the $p_b$ mean moments
and the $p_b$ intercept parameters cancel.
\end{proposition}

\begin{proof}
$\partial\mu_g/\partial\nu_g = I$ and, with $\alpha_g = 0$,
$\partial\mu_g/\partial\Lambda_g = 0$ and $\partial\mu_g/\partial(\text{cov}) =
0$, so the mean rows of $\Delta^{+}$ are $(I, 0)$. The intercept rows of $J^{+}$
are $(I, 0)$ and the covariance-parameter rows are $(0, J_b)$, since perturbing
$S_b$ leaves $\nu_g = m_g$ fixed. Multiplying gives the block-diagonal
$\Delta^{+}J^{+}$ above, and the zero top-left block of $M^{+}$ annihilates the
mean moments in $M^{+\mathsf T} V^{+} M^{+}$.
\end{proof}

Proposition~\ref{prop:mean-inert} means only the standard errors change: the
augmented $\Om$ picks up the intercept block $\Sigma_b/N_b$ (and, under a
distribution-free $\Gamma$, an intercept-by-loading cross-block from the third
moments). Intercept equality with $\alpha = 0$ is then just another
$R\theta = c$ in \eqref{eq:proj}.

\section{True scalar invariance: the reference-group mean map}
\label{sec:scalar}

Scalar invariance frees the latent means. Fix a reference group $r$ with
$\alpha_r = 0$. Given the common loadings read off the reference block,
$\Lambda \equiv \Lambda_r$, the closed-form map solves for the intercepts and
the latent means by least squares through $\operatorname{col}(\Lambda)$:
\begin{equation}
  \nu = m_r,
  \qquad
  \alpha_g = (\Lambda\T\Lambda)^{-1}\Lambda\T(m_g - m_r),
  \qquad
  d_g = (I - P_\Lambda)(m_g - m_r),
  \label{eq:refmap}
\end{equation}
with $P_\Lambda = \Lambda(\Lambda\T\Lambda)^{-1}\Lambda\T$. Scalar invariance
holds if and only if every $d_g = 0$, that is, the group mean difference lies in
$\operatorname{col}(\Lambda)$.

The residual $d_g$ is bilinear in $(\Lambda, m)$, so the test is a linearized
Wald rather than a projection. To leading order, using $r = d_g \approx 0$ under
the null,
\begin{equation}
  \partial d_g \approx -(I - P_\Lambda)\,(\partial\Lambda)\,\alpha_g
                       + (I - P_\Lambda)\,(\partial m_g - \partial m_r),
  \label{eq:dg-jac}
\end{equation}
which gives the stacked Jacobian $G = \partial d/\partial\theta$ and
\begin{equation}
  W = d\T \Cov(d)^{+} d,
  \qquad
  \Cov(d) = G\,\Om\, G\T,
  \qquad
  \mathrm{df} = (G-1)(p - m),
  \label{eq:scalar-wald}
\end{equation}
where $m$ is the number of factors and $\Cov(d)$ is singular of rank
$(G-1)(p-m)$, so the statistic is a pseudo-inverse Wald.

\begin{proposition}[Rank and reference]
\label{prop:scalar-rank}
Each column of $G$ restricted to group $g$ lies in
$\operatorname{col}(I - P_\Lambda)$, and so does $d_g$. Consequently
$\rank\Cov(d) = (G-1)(p-m)$ generically and $d$ lies in the range of $\Cov(d)$
almost surely, so $W \to \chi^2_{(G-1)(p-m)}$. The non-reference groups are
coupled through the shared reference mean $m_r$ and loadings $\Lambda$, and
$\Om$ supplies that coupling automatically through \eqref{eq:omega-blockdiag}.
\end{proposition}

\begin{remark}
The map \eqref{eq:refmap} uses the reference group's configural loadings as the
common metric, which is the natural choice for a reference-group parameterization
and keeps $\Cov(d)$ exact through block diagonality. A pooled-loading refinement
(test given metric invariance) would use the metric-projected common loadings of
\eqref{eq:proj}; it is more efficient but couples the mean and loading blocks,
and is left as future work.
\end{remark}

\section{Novelty and prior art}

The single-group closed-form CFA point estimator is Guttman (1952), and its
single-factor reliability specialization is Hancock and An (2020). Dhaene and
Rosseel (2024) study closed-form and structured estimators in the same spirit.
The residual-based delta-method inference bundle is developed in the companion
notes. What appears new here is the combination: a multiple-group closed-form
CFA with a block-diagonal delta-method covariance, an $\Om^{-1}$ minimum-distance
projection that turns any linear invariance hypothesis into an exact $\chi^2$
test through Proposition~\ref{prop:duality}, and a closed-form reference-group
scalar map with a linearized invariance Wald
(Proposition~\ref{prop:scalar-rank}). A full external prior-art check remains to
be completed before any publication claim.

\begin{thebibliography}{9}

\bibitem{guttman1952}
L.~Guttman.
\newblock Multiple group methods for common-factor analysis: their basis,
  computation, and interpretation.
\newblock \emph{Psychometrika}, 17(2):209--222, 1952.
\newblock \href{https://doi.org/10.1007/BF02288782}{doi:10.1007/BF02288782}.

\bibitem{hancockan2020}
G.~R. Hancock and J.~An.
\newblock A closed-form alternative for estimating $\omega$ reliability under
  unidimensionality.
\newblock \emph{Measurement: Interdisciplinary Research and Perspectives},
  18(1):1--14, 2020.
\newblock \href{https://doi.org/10.1080/15366367.2019.1656049}{doi:10.1080/15366367.2019.1656049}.

\bibitem{dhaenerosseel2024}
S.~Dhaene and Y.~Rosseel.
\newblock Structured and closed-form estimators in structural equation modeling.
\newblock Working paper, 2024.

\end{thebibliography}

\end{document}
